% --- Archon physics formalization source begin ---
% archon:physics
% archon:covers PhyXMiniProblems/problem_phyx_mini_0521.lean
% archon:source-report reports/phyx_mini/problem_phyx_mini_0521.source.json

\chapter{Physics problem phyx\_mini\_0521}
\label{ch:PhyXMiniProblems_problem_phyx_mini_0521}

\paragraph{Problem source.}
\#\# Physical scenario

Consider a potential well defined as \textbackslash{}( U(x) = \textbackslash{}infty \textbackslash{}) for \textbackslash{}( x < 0 \textbackslash{}), \textbackslash{}( U(x) = 0 \textbackslash{}) for \textbackslash{}( 0 < x < L \textbackslash{}), and \textbackslash{}( U(x) = U\_0 > 0 \textbackslash{}) for \textbackslash{}( x > L \textbackslash{}) as shown in figure. Consider a particle with mass \textbackslash{}( m \textbackslash{}) and kinetic energy \textbackslash{}( E < U\_0 \textbackslash{}) that is trapped in the well.

\#\# Question

The wave function must remain finite as \textbackslash{}( x \textbackslash{}to \textbackslash{}infty \textbackslash{}). What must the form of the function \textbackslash{}( \textbackslash{}psi(x) \textbackslash{}) for \textbackslash{}( x > L \textbackslash{}) be in order to satisfy both the Schrödinger equation and this boundary condition at infinity?

\#\# Answer choices

- A: A: \textbackslash{}( \textbackslash{}kappa\textasciicircum{}2 = m(U\_0 + E)/\textbackslash{}hbar

- B: B: \textbackslash{}( \textbackslash{}kappa\textasciicircum{}2 = m(U\_0 - E)/\textbackslash{}hbar

- C: C: \textbackslash{}( \textbackslash{}kappa\textasciicircum{}2 = 2m(U\_0 + E)/\textbackslash{}hbar

- D: D: \textbackslash{}\textbackslash{}( \textbackslash{}kappa\textasciicircum{}2 = 2m(U\_0 - E)/\textbackslash{}hbar

\#\# Figure caption (auxiliary; use the image as primary evidence)

The image is a graph illustrating a potential well, labeled with axes \textbackslash{}( U(x) \textbackslash{}) (vertical) and \textbackslash{}( x \textbackslash{}) (horizontal). Here's a detailed description:

- The vertical axis \textbackslash{}( U(x) \textbackslash{}) represents potential energy.
- The horizontal axis \textbackslash{}( x \textbackslash{}) represents position.
- At \textbackslash{}( x = 0 \textbackslash{}), the potential energy is infinite, indicated by an upward arrow.
- From \textbackslash{}( x = 0 \textbackslash{}) to \textbackslash{}( x = L \textbackslash{}), there is a constant potential denoted by \textbackslash{}( U\_0 \textbackslash{}). It is marked by a horizontal line extending to \textbackslash{}( x = L \textbackslash{}).
- Beyond \textbackslash{}( x = L \textbackslash{}), the potential drops to zero, depicted by a horizontal line at the bottom.
- There is a dashed horizontal line across from \textbackslash{}( 0 \textbackslash{}) to \textbackslash{}( L \textbackslash{}) at the level of \textbackslash{}( U\_0 \textbackslash{}).
- The region from \textbackslash{}( 0 \textbackslash{}) to \textbackslash{}( L \textbackslash{}) is shaded, highlighting the width of the potential well.

This graph is a typical representation of a finite potential well in quantum mechanics.

\#\# Dataset metadata

Quantum Phenomena; Physical Model Grounding Reasoning, Numerical Reasoning

\paragraph{Recorded answer/context.}
D: D: \textbackslash{}\textbackslash{}( \textbackslash{}kappa\textasciicircum{}2 = 2m(U\_0 - E)/\textbackslash{}hbar

\paragraph{Figure/image path.}
/root/proposal\_for\_physic/science-mango/phyx\_mini\_run/phyx\_data/test\_image/521.png

\paragraph{Formalization target.}
create a compiling Lean file with sorry bodies at `PhyXMiniProblems/problem\_phyx\_mini\_0521.lean`. The Lean declarations must preserve the physical quantities, dimensions or dimensional roles, figure labels, governing-law hypotheses, and final relation expressed by this problem.
Use Mathlib/Physlib names found through LeanExplore where available. If a physics API is missing, introduce faithful local abstractions rather than scalar placeholder aliases.

\begin{theorem}[Physics formalization target]
\label{thm:physics:phyx_mini_0521:target}
\lean{PhyXMiniProblems.ProblemPhyXMini0521.exterior_waveFunction_has_decaying_exponential_form}
\uses{def:physics:phyx-mini-0521:phyxminiproblems-problemphyxmini0521-wavenumberquantity, def:physics:phyx-mini-0521:phyxminiproblems-problemphyxmini0521-lengthinmeters, def:physics:phyx-mini-0521:phyxminiproblems-problemphyxmini0521-massinkilograms, def:physics:phyx-mini-0521:phyxminiproblems-problemphyxmini0521-energyinjoules, def:physics:phyx-mini-0521:phyxminiproblems-problemphyxmini0521-actioninjouleseconds, def:physics:phyx-mini-0521:phyxminiproblems-problemphyxmini0521-wavenumberininversemeters, def:physics:phyx-mini-0521:phyxminiproblems-problemphyxmini0521-finitestepwellsetup, def:physics:phyx-mini-0521:phyxminiproblems-problemphyxmini0521-matchesfinitestepwellscenario, def:physics:phyx-mini-0521:phyxminiproblems-problemphyxmini0521-matchessuppliedstepwellfigure, def:physics:phyx-mini-0521:phyxminiproblems-problemphyxmini0521-hasphysicalfinitestepwellparameters, def:physics:phyx-mini-0521:phyxminiproblems-problemphyxmini0521-usesstandardreducedplanckconstant, def:physics:phyx-mini-0521:phyxminiproblems-problemphyxmini0521-satisfiesexteriorstationaryschrodingerequation, def:physics:phyx-mini-0521:phyxminiproblems-problemphyxmini0521-satisfiesfiniteboundaryconditionatinfinity}
The assigned autoformalize agent should translate this physics problem into Lean declarations in the covered file, with theorem and lemma proof bodies written as `by sorry`.
\end{theorem}
\begin{proof}
This is an autoformalization task, not a proof task. Produce statements that can later be proved by the physics prover without weakening the physical model.
\end{proof}

% --- Archon declaration topology begin ---
\section{Declaration topology}
\begin{definition}[Length Quantity]
  \label{def:physics:phyx-mini-0521:phyxminiproblems-problemphyxmini0521-lengthquantity}
  \lean{PhyXMiniProblems.ProblemPhyXMini0521.LengthQuantity}
  A nonnegative, unit-independent physical length
\end{definition}
\begin{proof}
  This is the declared model definition of Length Quantity.
\end{proof}
\begin{definition}[Mass Quantity]
  \label{def:physics:phyx-mini-0521:phyxminiproblems-problemphyxmini0521-massquantity}
  \lean{PhyXMiniProblems.ProblemPhyXMini0521.MassQuantity}
  A nonnegative, unit-independent physical mass
\end{definition}
\begin{proof}
  This is the declared model definition of Mass Quantity.
\end{proof}
\begin{definition}[action Dimension]
  \label{def:physics:phyx-mini-0521:phyxminiproblems-problemphyxmini0521-actiondimension}
  \lean{PhyXMiniProblems.ProblemPhyXMini0521.actionDimension}
  The physical dimension of action, mass * length² / time
\end{definition}
\begin{proof}
  This is the declared model definition of action Dimension.
\end{proof}
\begin{definition}[Action Quantity]
  \label{def:physics:phyx-mini-0521:phyxminiproblems-problemphyxmini0521-actionquantity}
  \lean{PhyXMiniProblems.ProblemPhyXMini0521.ActionQuantity}
  \uses{def:physics:phyx-mini-0521:phyxminiproblems-problemphyxmini0521-actiondimension}
  A nonnegative, unit-independent physical action
\end{definition}
\begin{proof}
  This is the declared model definition of Action Quantity.
\end{proof}
\begin{definition}[Wave Number Quantity]
  \label{def:physics:phyx-mini-0521:phyxminiproblems-problemphyxmini0521-wavenumberquantity}
  \lean{PhyXMiniProblems.ProblemPhyXMini0521.WaveNumberQuantity}
  A nonnegative, unit-independent inverse length (wave number)
\end{definition}
\begin{proof}
  This is the declared model definition of Wave Number Quantity.
\end{proof}
\begin{definition}[length In Meters]
  \label{def:physics:phyx-mini-0521:phyxminiproblems-problemphyxmini0521-lengthinmeters}
  \lean{PhyXMiniProblems.ProblemPhyXMini0521.lengthInMeters}
  \uses{def:physics:phyx-mini-0521:phyxminiproblems-problemphyxmini0521-lengthquantity}
  Metre readout of a physical length
\end{definition}
\begin{proof}
  This is the declared model definition of length In Meters.
\end{proof}
\begin{definition}[mass In Kilograms]
  \label{def:physics:phyx-mini-0521:phyxminiproblems-problemphyxmini0521-massinkilograms}
  \lean{PhyXMiniProblems.ProblemPhyXMini0521.massInKilograms}
  \uses{def:physics:phyx-mini-0521:phyxminiproblems-problemphyxmini0521-massquantity}
  Kilogram readout of a physical mass
\end{definition}
\begin{proof}
  This is the declared model definition of mass In Kilograms.
\end{proof}
\begin{definition}[energy In Joules]
  \label{def:physics:phyx-mini-0521:phyxminiproblems-problemphyxmini0521-energyinjoules}
  \lean{PhyXMiniProblems.ProblemPhyXMini0521.energyInJoules}
  Joule readout of a physical energy
\end{definition}
\begin{proof}
  This is the declared model definition of energy In Joules.
\end{proof}
\begin{definition}[action In Joule Seconds]
  \label{def:physics:phyx-mini-0521:phyxminiproblems-problemphyxmini0521-actioninjouleseconds}
  \lean{PhyXMiniProblems.ProblemPhyXMini0521.actionInJouleSeconds}
  \uses{def:physics:phyx-mini-0521:phyxminiproblems-problemphyxmini0521-actionquantity}
  Joule-second readout of a physical action
\end{definition}
\begin{proof}
  This is the declared model definition of action In Joule Seconds.
\end{proof}
\begin{definition}[wave Number In Inverse Meters]
  \label{def:physics:phyx-mini-0521:phyxminiproblems-problemphyxmini0521-wavenumberininversemeters}
  \lean{PhyXMiniProblems.ProblemPhyXMini0521.waveNumberInInverseMeters}
  \uses{def:physics:phyx-mini-0521:phyxminiproblems-problemphyxmini0521-wavenumberquantity}
  Inverse-metre readout of a physical wave number
\end{definition}
\begin{proof}
  This is the declared model definition of wave Number In Inverse Meters.
\end{proof}
\begin{definition}[zero Potential Energy]
  \label{def:physics:phyx-mini-0521:phyxminiproblems-problemphyxmini0521-zeropotentialenergy}
  \lean{PhyXMiniProblems.ProblemPhyXMini0521.zeroPotentialEnergy}
  The unit-independent physical zero of energy
\end{definition}
\begin{proof}
  This is the declared model definition of zero Potential Energy.
\end{proof}
\begin{definition}[Potential Value]
  \label{def:physics:phyx-mini-0521:phyxminiproblems-problemphyxmini0521-potentialvalue}
  \lean{PhyXMiniProblems.ProblemPhyXMini0521.PotentialValue}
  A potential value is either a finite physical energy or an infinite wall
\end{definition}
\begin{proof}
  This is the declared model definition of Potential Value.
\end{proof}
\begin{definition}[Figure Label]
  \label{def:physics:phyx-mini-0521:phyxminiproblems-problemphyxmini0521-figurelabel}
  \lean{PhyXMiniProblems.ProblemPhyXMini0521.FigureLabel}
  Literal labels visible in the supplied potential-energy graph
\end{definition}
\begin{proof}
  This is the declared model definition of Figure Label.
\end{proof}
\begin{definition}[Supplied Step Well Figure]
  \label{def:physics:phyx-mini-0521:phyxminiproblems-problemphyxmini0521-suppliedstepwellfigure}
  \lean{PhyXMiniProblems.ProblemPhyXMini0521.SuppliedStepWellFigure}
  \uses{def:physics:phyx-mini-0521:phyxminiproblems-problemphyxmini0521-figurelabel}
  Qualitative and textual facts read from the primary raster
\end{definition}
\begin{proof}
  This is the declared model definition of Supplied Step Well Figure.
\end{proof}
\begin{definition}[Finite Step Well Setup]
  \label{def:physics:phyx-mini-0521:phyxminiproblems-problemphyxmini0521-finitestepwellsetup}
  \lean{PhyXMiniProblems.ProblemPhyXMini0521.FiniteStepWellSetup}
  \uses{def:physics:phyx-mini-0521:phyxminiproblems-problemphyxmini0521-lengthquantity, def:physics:phyx-mini-0521:phyxminiproblems-problemphyxmini0521-massquantity, def:physics:phyx-mini-0521:phyxminiproblems-problemphyxmini0521-actionquantity, def:physics:phyx-mini-0521:phyxminiproblems-problemphyxmini0521-potentialvalue, def:physics:phyx-mini-0521:phyxminiproblems-problemphyxmini0521-suppliedstepwellfigure}
  The independent physical quantities in the problem. The coordinate argument of potentialAtMeterCoordinate and waveFunctionAmplitude is a metre readout. The energy called E in the source is kinetic energy in the zero-potential interior, hence also the stationary state's total-energy readout
\end{definition}
\begin{proof}
  This is the declared model definition of Finite Step Well Setup.
\end{proof}
\begin{definition}[Matches Finite Step Well Scenario]
  \label{def:physics:phyx-mini-0521:phyxminiproblems-problemphyxmini0521-matchesfinitestepwellscenario}
  \lean{PhyXMiniProblems.ProblemPhyXMini0521.MatchesFiniteStepWellScenario}
  \uses{def:physics:phyx-mini-0521:phyxminiproblems-problemphyxmini0521-lengthinmeters, def:physics:phyx-mini-0521:phyxminiproblems-problemphyxmini0521-zeropotentialenergy, def:physics:phyx-mini-0521:phyxminiproblems-problemphyxmini0521-finitestepwellsetup}
  The piecewise potential and trapped-state description stated in the prose and confirmed by the primary image. No exterior wave-function form or decay rate occurs in this premise
\end{definition}
\begin{proof}
  This is the declared model definition of Matches Finite Step Well Scenario.
\end{proof}
\begin{definition}[Matches Supplied Step Well Figure]
  \label{def:physics:phyx-mini-0521:phyxminiproblems-problemphyxmini0521-matchessuppliedstepwellfigure}
  \lean{PhyXMiniProblems.ProblemPhyXMini0521.MatchesSuppliedStepWellFigure}
  \uses{def:physics:phyx-mini-0521:phyxminiproblems-problemphyxmini0521-finitestepwellsetup}
  Every label and qualitative feature visible in image 521
\end{definition}
\begin{proof}
  This is the declared model definition of Matches Supplied Step Well Figure.
\end{proof}
\begin{definition}[Has Physical Finite Step Well Parameters]
  \label{def:physics:phyx-mini-0521:phyxminiproblems-problemphyxmini0521-hasphysicalfinitestepwellparameters}
  \lean{PhyXMiniProblems.ProblemPhyXMini0521.HasPhysicalFiniteStepWellParameters}
  \uses{def:physics:phyx-mini-0521:phyxminiproblems-problemphyxmini0521-lengthinmeters, def:physics:phyx-mini-0521:phyxminiproblems-problemphyxmini0521-massinkilograms, def:physics:phyx-mini-0521:phyxminiproblems-problemphyxmini0521-energyinjoules, def:physics:phyx-mini-0521:phyxminiproblems-problemphyxmini0521-actioninjouleseconds, def:physics:phyx-mini-0521:phyxminiproblems-problemphyxmini0521-finitestepwellsetup}
  Positivity and ordering conditions on the physical parameters. In particular, E < U₀ is source data, not the requested formula for κ
\end{definition}
\begin{proof}
  This is the declared model definition of Has Physical Finite Step Well Parameters.
\end{proof}
\begin{definition}[Uses Standard Reduced Planck Constant]
  \label{def:physics:phyx-mini-0521:phyxminiproblems-problemphyxmini0521-usesstandardreducedplanckconstant}
  \lean{PhyXMiniProblems.ProblemPhyXMini0521.UsesStandardReducedPlanckConstant}
  \uses{def:physics:phyx-mini-0521:phyxminiproblems-problemphyxmini0521-actioninjouleseconds, def:physics:phyx-mini-0521:phyxminiproblems-problemphyxmini0521-finitestepwellsetup}
  Physlib's Constants.ℏ is the standard reduced Planck constant in joule-seconds. This calibration does not constrain the requested decay rate
\end{definition}
\begin{proof}
  This is the declared model definition of Uses Standard Reduced Planck Constant.
\end{proof}
\begin{definition}[exterior Schrodinger Left Hand Side]
  \label{def:physics:phyx-mini-0521:phyxminiproblems-problemphyxmini0521-exteriorschrodingerlefthandside}
  \lean{PhyXMiniProblems.ProblemPhyXMini0521.exteriorSchrodingerLeftHandSide}
  \uses{def:physics:phyx-mini-0521:phyxminiproblems-problemphyxmini0521-massinkilograms, def:physics:phyx-mini-0521:phyxminiproblems-problemphyxmini0521-energyinjoules, def:physics:phyx-mini-0521:phyxminiproblems-problemphyxmini0521-actioninjouleseconds, def:physics:phyx-mini-0521:phyxminiproblems-problemphyxmini0521-finitestepwellsetup}
  The kinetic-plus-potential side of the one-dimensional stationary Schrödinger equation in the exterior constant-potential region
\end{definition}
\begin{proof}
  This is the declared model definition of exterior Schrodinger Left Hand Side.
\end{proof}
\begin{definition}[Satisfies Exterior Stationary Schrodinger Equation]
  \label{def:physics:phyx-mini-0521:phyxminiproblems-problemphyxmini0521-satisfiesexteriorstationaryschrodingerequation}
  \lean{PhyXMiniProblems.ProblemPhyXMini0521.SatisfiesExteriorStationarySchrodingerEquation}
  \uses{def:physics:phyx-mini-0521:phyxminiproblems-problemphyxmini0521-lengthinmeters, def:physics:phyx-mini-0521:phyxminiproblems-problemphyxmini0521-energyinjoules, def:physics:phyx-mini-0521:phyxminiproblems-problemphyxmini0521-finitestepwellsetup, def:physics:phyx-mini-0521:phyxminiproblems-problemphyxmini0521-exteriorschrodingerlefthandside}
  Twice-continuous differentiability and the stationary Schrödinger equation on x > L. This is the governing law; it does not state an exponential ansatz or a formula for its decay constant
\end{definition}
\begin{proof}
  This is the declared model definition of Satisfies Exterior Stationary Schrodinger Equation.
\end{proof}
\begin{definition}[Satisfies Finite Boundary Condition At Infinity]
  \label{def:physics:phyx-mini-0521:phyxminiproblems-problemphyxmini0521-satisfiesfiniteboundaryconditionatinfinity}
  \lean{PhyXMiniProblems.ProblemPhyXMini0521.SatisfiesFiniteBoundaryConditionAtInfinity}
  \uses{def:physics:phyx-mini-0521:phyxminiproblems-problemphyxmini0521-lengthinmeters, def:physics:phyx-mini-0521:phyxminiproblems-problemphyxmini0521-finitestepwellsetup}
  The wave function remains finite at positive infinity: beyond some exterior cutoff its norm has a finite uniform bound. This excludes a growing mode but does not assume the requested decaying-exponential form or a zero limit
\end{definition}
\begin{proof}
  This is the declared model definition of Satisfies Finite Boundary Condition At Infinity.
\end{proof}
% --- Archon declaration topology end ---
% --- Archon physics formalization source end ---
